\section{Absorbable directions}\label{sec:absorbable}

The rotation of Section~\ref{sec:rotation} is exact for any fixed
orthogonal matrix; whether it helps depends on the basis it is built from.
The basis should contain every direction the fixed effects can absorb,
and only such directions. This section characterizes them.

Fix a random-effect column $c$ and write $z_i = Z_i e_c \in \R^{n_i}$ for
subject $i$'s entries of it; in this section only, $z_i(t)$ denotes the
scalar entry of column $c$ in the design row at time $t$, rather than the
whole row of (\ref{eq:traj}).

\begin{definition}[Absorbable direction]\label{def:abs}
A vector $v \in \R^N$ is absorbable for column $c$ if there is a witness
$d \in \R^p$ such that the shift $B e_c \mapsto B e_c + v$, $\beta \mapsto
\beta + d$ leaves every linear predictor and every marker trajectory, and
hence the likelihood, unchanged:
\begin{align}
X_i d + v_i z_i &= 0 \quad \text{for every } i, \tag{L}\label{eq:L}\\
x_i(t)^{\top} d + v_i\, z_i(t) &= 0 \quad \text{for every } i
\text{ and every } t \in [0, \tau]. \tag{T}\label{eq:T}
\end{align}
Write $\A_c$ for the set of absorbable directions and $\A_c^{\mathrm{long}}$
for the set satisfying (\ref{eq:L}) alone, which is the whole requirement
in a mixed model without a survival submodel.
\end{definition}

Both sets are linear subspaces of $\R^N$. The likelihood uses $m_i(t)$
only up to each subject's event or censoring time, so requiring
(\ref{eq:T}) on all of $[0, \tau]$ is slightly more than likelihood
invariance needs; it makes $\A_c$ a structural property of the model
rather than of the observed follow-up. It is sufficient for likelihood
invariance, not necessarily the largest set of shifts that leaves a
particular observed likelihood unchanged. The prior on the random effects is
not invariant under these shifts, and that asymmetry is what identifies
the split at all.

\subsection{Characterization and construction}

Let $\T: \R^p \to \bigoplus_i \R^{n_i}$, $(\T d)_i = X_i d$, map a shift
of $\beta$ to its effect on the stacked linear predictor, and $\U_c: \R^N
\to \bigoplus_i \R^{n_i}$, $(\U_c v)_i = v_i z_i$, do the same for a shift
of random-effect column $c$. Let $P_{\T}$ be the orthogonal projector onto
$\range \T$ and, for $z_i \neq 0$, $P_{z_i}$ the projector onto
$\spn(z_i)$. Let $S_c = \{i : z_i \neq 0\}$ be the subjects with at least
one observation at which column $c$ is nonzero, and let $M_c \in
\R^{(\sum_i n_i) \times p}$ stack $(I - P_{z_i}) X_i$ for $i \in S_c$ and
$X_i$ for $i \notin S_c$.

\begin{proposition}[Characterization]\label{prop:char}
\begin{enumerate}
\item[(a)] $\A_c^{\mathrm{long}} = \{v : \U_c v \in \range \T\}$, the
preimage of $\range \T$ under $\U_c$, and equals $\ker\{(I -
P_{\T})\,\U_c\}$. It depends on the fixed-effects design only
through its column space.
\item[(b)] A vector $d$ witnesses some $v \in \A_c^{\mathrm{long}}$ if and
only if $d \in \ker M_c$. The witnessed $v$ then has $v_i = -z_i^{\top}
X_i d / z_i^{\top} z_i$ for $i \in S_c$, and $v_i$ is unrestricted for
$i \notin S_c$.
\end{enumerate}
\end{proposition}

\begin{proof}
(a) Condition (\ref{eq:L}) says $\T d = -\U_c v$ for some $d$, that is,
$\U_c v \in \range \T$, which is equivalent to $(I - P_{\T})\U_c v = 0$.
(b) For $i \in S_c$, $X_i d = -v_i z_i$ holds for some scalar $v_i$ if and
only if $X_i d \in \spn(z_i)$, that is, $(I - P_{z_i}) X_i d = 0$, and then
$v_i = -z_i^{\top} X_i d / z_i^{\top} z_i$. For $i \notin S_c$,
(\ref{eq:L}) reads $X_i d = 0$ whatever $v_i$ is.
\end{proof}

Part (b) is the construction used in practice. In a mixed model without
a survival submodel, a subject with $z_i = 0$ has no random effect in
column $c$ in the likelihood, so we may take $S_c$ to be every subject.
One singular value decomposition of the $(\sum_i n_i) \times
p$ matrix $M_c$ gives a basis of $\ker M_c$, each basis vector $d$ gives one
direction $v(d)$, and, if $X$ has full column rank, $d \mapsto v(d)$ is
injective, so $\dim \A_c^{\mathrm{long}} = \dim \ker M_c \leq p$. The
computation is done once, before sampling. In the
implementation every returned direction is also checked against
(\ref{eq:L}) before use, so a loose rank tolerance cannot admit a direction
that is not absorbable.

\subsection{The survival submodel}

In a joint model the requirement (\ref{eq:T}) quantifies over all
$t \in [0, \tau]$, while $M_c$ sees only the observation times.

\begin{condition}[Structural extension]\label{cond:S}
For every $d \in \ker M_c$ and every subject $i$ there is a scalar
$v_i(d)$, equal to $-z_i^{\top} X_i d / z_i^{\top} z_i$ when $i \in S_c$,
such that $x_i(t)^{\top} d + v_i(d)\, z_i(t) = 0$ for every $t \in [0,
\tau]$. (Under the assumption of Proposition~\ref{prop:joint}, $v_i(d)$ is
then unique for every $i$.)
\end{condition}

\begin{proposition}[Joint model]\label{prop:joint}
Suppose every subject $i \notin S_c$ has some $t \in [0, \tau]$ with
$z_i(t) \neq 0$. Under Condition~\ref{cond:S}, $\A_c = \{v(d) : d \in \ker
M_c\}$. For a subject $i \notin S_c$, $v_i(d)$ is determined by
(\ref{eq:T}) at any $t$ with $z_i(t) \neq 0$; in particular it is not
free, although the longitudinal data alone leave it free.
\end{proposition}

The proof is in Appendix~\ref{app:proofs}. Condition~\ref{cond:S}
holds for subject-constant columns (absorbed by the intercept column
$z_i(t) \equiv 1$), the time column (absorbed by a random slope $z_i(t) =
t$) and covariate-by-time interactions: in each, the identity holds as a
function of $t$, not by coincidence at the observed times. More generally
it holds when the functions $t \mapsto x_i(t)^{\top} d$ and $t \mapsto
z_i(t)$ are known to lie in a finite-dimensional function space
determined by its values at the subject's
distinct observation times, such as polynomials of degree below $n_i$
(Appendix~\ref{app:proofs}). For spline bases it can fail for subjects
with few observations relative to the flexibility of the basis. It is
cheap to check: \texttt{jmjax} evaluates the identity at the event times
and quadrature nodes at every fit and warns if the worst relative residual
exceeds $10^{-6}$. Proposition~\ref{prop:joint} also settles the entry of
a subject with no observation at which the column is nonzero, such as a
baseline-only subject for the slope column: it is forced by the survival
submodel, and setting it to zero would produce a direction the likelihood
does identify.

\subsection{A design where the natural check fails}

A natural construction, and the one \texttt{jmjax} originally used, looks
for columns of $X$ that are \emph{individually} subject-constant multiples
of $z$: column $j$ is kept if $X_i e_j = s_i z_i$ for every $i$, and the
basis is the span of the resulting vectors $s$. This finds a subspace of
$\A_c^{\mathrm{long}}$, but not always all of it.

\begin{example}\label{ex:spline}
Let every subject be observed at the same $m \geq 3$ distinct times, let
$Z_i = [\,1,\ t\,]$ with $c$ the slope column, and let $X_i = [\,1,\ t +
\tfrac12 t^2,\ t^2\,]$. The ratios of the three columns to $t$ are $1/t$,
$1 + t/2$ and $t$, none constant, so the column-by-column search returns
nothing. Yet $d = (0, 1, -\tfrac12)$ gives $X_i d = t$, so $v =
-\mathbf{1}_N$ satisfies (\ref{eq:L}) and (\ref{eq:T}), and $\A_c
\neq \{0\}$.
\end{example}

The column space in Example~\ref{ex:spline} is $\spn\{1, t, t^2\}$,
what an orthogonal-polynomial or natural-spline mean in time produces (the
monomial spelling \texttt{time + I(time\^{}2)} is caught, because the
time column appears on its own). With such a mean and a random slope, the
column-by-column search silently leaves the slope degeneracy in place. By
Proposition~\ref{prop:char}(a) the absorbable set depends only on the
column space, and the null-space construction of
Proposition~\ref{prop:char}(b), which \texttt{jmjax} now uses, does not
depend on how that space is spanned.

Finally, let $Q_c$ be an orthonormal basis of $\A_c$ and $Q \in \R^{N
\times k}$ an orthonormal basis of $\spn\bigcup_c \A_c$. In the simulation
design of Section~\ref{sec:simulation} (an intercept, time, age and sex,
with a random intercept and slope) $\A_0 = \spn\{\mathbf{1}, \mathrm{age},
\mathrm{sex}\}$ and $\A_1 = \spn\{\mathbf{1}\}$, so $k = 3$.
